%!TEX root = ../main.tex
\section{Proof of Theorem \ref{thm:lqr-lift} }
\label{sec:app-proof-lqg-lift}

\begin{proof}
    From
      \[
      \bar s=\mathcal F_S s_0+\mathcal F_E \bar\varepsilon,\qquad
      \bar s^+=\mathcal F_S^+ s_0+\mathcal F_E^+ \bar\varepsilon,\qquad
      \bar a=\mathcal K_S s_0+\mathcal K_E \bar\varepsilon,
      \]
    and using $s_{t+1} = F s_t + B \varepsilon_t$, we write out the explicit forms:
      
    \medskip
    \noindent\textbf{State maps.}
    \begin{align}\label{eq:lqr-lift-state-maps}
      \mathcal F_S=\begin{bmatrix}
      I\\[2pt] F\\[2pt] F^2\\[-2pt]\vdots\\[2pt] F^{T-1}
      \end{bmatrix}\!\in\mathbb R^{nT\times n},
      \qquad
      \mathcal F_S^+=\begin{bmatrix}
      F\\[2pt] F^2\\[-2pt]\vdots\\[2pt] F^{T}
      \end{bmatrix}\!\in\mathbb R^{nT\times n},
    \end{align}
  
    \[
    \mathcal F_E=
    \begin{bmatrix}
    \;0\;&\;0\;&\cdots&\;0&\;0\\
    B&0&\cdots&0&0\\
    FB&B&\ddots&\vdots&0\\
    \vdots&\ddots&\ddots&0&0\\
    F^{T-2}B&\cdots&FB&B&0 
    \end{bmatrix}\!\in\mathbb R^{nT\times mT},
    \qquad
    \mathcal F_E^+=
    \begin{bmatrix}
    B&0&\cdots&0\\
    FB&B&\ddots&\vdots\\
    \vdots&\ddots&\ddots&0\\
    F^{T-1}B&\cdots&FB&B
    \end{bmatrix}\!\in\mathbb R^{nT\times mT}.
    \]
    Equivalently, for \(t,i\in\{1,\ldots,T\}\),
    \[
    [\mathcal F_E]_{t,i}=
    \begin{cases}
    F^{\,t-1-i}B,& i\le t-1,\\[2pt]
    0,& i\ge t,
    \end{cases}
    \qquad
    [\mathcal F_E^+]_{t,i}=
    \begin{cases}
    F^{\,t-i}B,& i\le t,\\[2pt]
    0,& i>t.
    \end{cases}
    \]
    
    \medskip
    \noindent\textbf{Action maps.}
  
    \[
    \mathcal K_S=
    \begin{bmatrix}
    K\\[2pt] KF\\[-2pt]\vdots\\[2pt] KF^{T-1}
    \end{bmatrix}\!\in\mathbb R^{mT\times n},
    \qquad
    \mathcal K_E=
    \begin{bmatrix}
    I_m&0&\cdots&0&0\\
    KB&I_m&\ddots&\vdots&0\\
    KFB&KB&\ddots&0&0\\
    \vdots&\ddots&\ddots&I_m&0\\
    KF^{T-2}B&\cdots&KFB& KB&I_m
    \end{bmatrix}\!\in\mathbb R^{mT\times mT}.
    \]
    Equivalently, for \(t,i\in\{1,\ldots,T\}\),
    \[
    [\mathcal K_E]_{t,i}=
    \begin{cases}
    K F^{\,t-1-i} B,& i\le t-1,\\[2pt]
    I_m,& i=t,\\[2pt]
    0,& i>t.
    \end{cases}
    \]
\end{proof}




%%%%%%%%% remark
%%%%%%%%%
%%%%%%%%%
%%%%%%%%%



% \section{Proof of Theorem \ref{thm:lqg-variance-T} }
% \label{sec:app-proof-lqg-variance-T}

% \begin{proof}
    
% Note that to use Theorem~\ref{thm:lqg-variance-one-step}, we need to separate $\bar s$ into $s_0$ and $\bar \varepsilon$ parts.
% \[
% \nabla_{\tilde K}\log\pi(\bar a\mid \bar s)=\overline\Sigma^{-1}\bar\varepsilon\,\bar s^\top = \overline\Sigma^{-1}\bar\varepsilon\, (s_0^\top \mathcal F_s^\top+\bar\varepsilon^\top \mathcal F_E^\top ),\qquad
% \]

% Let $\mathcal P$ be the linear map that
% keeps the $T$ diagonal $m\times n$ blocks of an $(mT)\times(nT)$ matrix and sums them:
% \[
% \mathcal P(M):=\sum_{t=0}^{T-1} M_{tt}\in\mathbb R^{m\times n},\qquad
% \|\mathcal P\|_{F\to F}\le \sqrt{T}.
% \]
% The shared-$K$ (original) REINFORCE estimator can be written as
% \[
% \widehat G_K
% =\mathcal P\!\Big((\overline\Sigma^{-1}\bar\varepsilon)\,\bar s^{\!\top}\Big)\,R_T
% =\underbrace{\mathcal P\!\Big((\overline\Sigma^{-1}\bar\varepsilon)s_0^{\!\top}\mathcal F_S^{\!\top}\Big)}_{=:A}\,R_T
% +\underbrace{\mathcal P\!\Big((\overline\Sigma^{-1}\bar\varepsilon)\bar\varepsilon^{\!\top}\mathcal F_E^{\!\top}\Big)}_{=:B}\,R_T .
% \]
% Thus $\widehat G_K=(A+B)\,R_T$ with
% \[
% A=\mathcal P\!\Big((\overline\Sigma^{-1}\bar\varepsilon)s_0^{\!\top}\mathcal F_S^{\!\top}\Big),\qquad
% B=\mathcal P\!\Big((\overline\Sigma^{-1}\bar\varepsilon)\bar\varepsilon^{\!\top}\mathcal F_E^{\!\top}\Big).
% \]


% Using $\|X+Y\|_{\mathrm{F}}^2\le 2\|X\|_{\mathrm{F}}^2+2\|Y\|_{\mathrm{F}}^2$.

% \[
% \Var_F(\widehat G_K) = \mathbb E[||\widehat G_K - \mathbb E [\widehat G_K]||_F^2] = \mathbb E[||AR_T - \mathbb E [AR_T] + BR_T - \mathbb E [BR_T]||_F^2]
% \;\le\;2\,\mathbb E\!\big[\|A\|_{\mathrm{F}}^2\,\tilde R_T^{\,2}\big]\;+\;2\,\mathbb E\!\big[\|B\|_{\mathrm{F}}^2\,\tilde R_T^{\,2}\big],
% \]

% Using $\|\mathcal P(MC)\|_{\mathrm{F}}\le \|\mathcal P\|_{F\to F}\,\|M\|_{\mathrm{F}}\,\|C\|_2$ gives
% \[
% \|A\|_{\mathrm{F}}\le \sqrt{T}\,\|\overline\Sigma^{-1}\bar\varepsilon\,s_0^{\!\top}\|_{\mathrm{F}}\,\|\mathcal F_S\|_2,\qquad
% \|B\|_{\mathrm{F}}\le \sqrt{T}\,\|\overline\Sigma^{-1}\bar\varepsilon\,\bar\varepsilon^{\!\top}\|_{\mathrm{F}}\,\|\mathcal F_E\|_2.
% \]
% Therefore
% \begin{equation}\label{eq:K-variance-split}
% \Var_F(\widehat G_K)
% \;\le\;2T\,\|\mathcal F_S\|_2^2\,\underbrace{\mathbb E\!\big[\|\overline\Sigma^{-1}\bar\varepsilon\,s_0^{\!\top}\|_{\mathrm{F}}^2\,R_T^{\,2}\big]}_{(\star)}
% \;+\;2T\,\|\mathcal F_E\|_2^2\,\underbrace{\mathbb E\!\big[\|\overline\Sigma^{-1}\bar\varepsilon\,\bar\varepsilon^{\!\top}\|_{\mathrm{F}}^2\,R_T^{\,2}\big]}_{(\dagger)} .
% \end{equation}

% The factor inside $(\star)$ is exactly of the one-step form with the lifted blocks
% \((\overline M_{ss},\overline M_{se},\overline M_{ee})\) from Theorem~\ref{thm:lqr-lift} and noise covariance \(\overline\Sigma\).
% Thus can be directly bounded by Theorem~\ref{thm:lqg-variance-one-step}.

% Term $(\dagger)$ can be bounded using the same technic as in Theorem~\ref{thm:lqg-variance-one-step}. To write everything down:

% \begin{align}
%   (\star) &= A_\star + B_\star + C_\star + D_\star,\quad\\
%   (\dagger) &= A_\dagger + B_\dagger + C_\dagger + D_\dagger,\\
%   A_\star
% &=\tr(\overline\Sigma^{-1})\Big[
%       \tr(\Sigma_0)\,\tr(\Sigma_0 \overline M_{ss})^2
%       +4\,\tr(\Sigma_0^2 \overline M_{ss})\,\tr(\Sigma_0 \overline M_{ss})
%       +2\,\tr(\Sigma_0 \overline M_{ss}\Sigma_0 \overline M_{ss})\,\tr(\Sigma_0) \notag \\
%       &~+8\,\tr(\Sigma_0^2 \overline M_{ss}\Sigma_0 \overline M_{ss})
%     \Big],\\[4pt]
% B_\star
% &=\tr(\Sigma_0)\Big[
%       8\,\tr(\overline\Sigma^{-1}\overline M_{ee}\,\overline\Sigma\,\overline M_{ee}\,\overline\Sigma)
%       +4\,\tr(\overline\Sigma^{-1}\overline M_{ee}\,\overline\Sigma)\,\tr(\overline\Sigma\,\overline M_{ee})
%       +2\,\tr(\overline\Sigma\,\overline M_{ee}\,\overline\Sigma\,\overline M_{ee})\,\tr(\overline\Sigma^{-1}) \notag \\
%       &~+\tr(\overline\Sigma^{-1})\,\tr(\overline\Sigma\,\overline M_{ee})^2
%     \Big],\\[4pt]
% C_\star
% &=2\Big[
%       2\,\tr(\Sigma_0^2 \overline M_{ss})+\tr(\Sigma_0)\tr(\Sigma_0 \overline M_{ss})
%     \Big]\Big[
%       2\,\tr(\overline\Sigma^{-1}\overline M_{ee}\,\overline\Sigma)+\tr(\overline\Sigma^{-1})\tr(\overline\Sigma\,\overline M_{ee})
%     \Big],\\[4pt]
% D_\star
% &=4\Big\{
%       2\big[\,2\,\tr(\Sigma_0^2 \overline U)+\tr(\Sigma_0)\tr(\Sigma_0 \overline U)\big]
%       +\tr(\overline\Sigma^{-1})\big[\,2\,\tr(\Sigma_0^2 \overline V)+\tr(\Sigma_0)\tr(\Sigma_0 \overline V)\big]
%     \Big\}.\\[6pt]
%     A_\dagger
%     &=\big(2\tr(I_{mT})+\tr(\overline\Sigma^{-1})\,\tr(\overline\Sigma)\big)
%        \Big[\,2\,\tr(\Sigma_0 \overline M_{ss}\Sigma_0 \overline M_{ss})+\tr(\Sigma_0 \overline M_{ss})^2\Big],\\[6pt]
%     C_\dagger
%     &=2\,\tr(\Sigma_0 \overline M_{ss})\Big[
%           8\,\tr(\overline M_{ee}\,\overline\Sigma)
%           +2\Big(\tr(I_{mT})\,\tr(\overline M_{ee}\,\overline\Sigma)
%                   +\tr(\overline\Sigma^{-1}\overline M_{ee}\,\overline\Sigma)\,\tr(\overline\Sigma) \notag \\
%                   &~+\tr(\overline\Sigma\,\overline M_{ee}\,\overline\Sigma)\,\tr(\overline\Sigma^{-1})\Big)
%           +\tr(\overline\Sigma^{-1})\,\tr(\overline\Sigma)\,\tr(\overline M_{ee}\,\overline\Sigma)
%        \Big],\\[6pt]
%     D_\dagger
%     &=4\Big\{
%           8\,\tr(\Sigma_0 \overline V)
%           +2\Big(\tr(I_{mT})\,\tr(\Sigma_0 \overline V)
%                   +\tr(\Sigma_0 \overline U)\,\tr(\overline\Sigma)
%                   +\tr(\Sigma_0 \overline W)\,\tr(\overline\Sigma^{-1})\Big) \notag \\
%           &+\tr(\overline\Sigma^{-1})\,\tr(\overline\Sigma)\,\tr(\Sigma_0 \overline V)
%        \Big\},\\[6pt]
%     B_\dagger
%     &=\Psi_4\!\big(\overline\Sigma^{-2},\,I_{mT},\,\overline M_{ee},\,\overline M_{ee}\big).
% \end{align}
% where 
% \[
% \Psi_4(A,B,C,D)
% :=\mathbb E\big[(\bar\varepsilon^\top A\bar\varepsilon)(\bar\varepsilon^\top B\bar\varepsilon)
%                 (\bar\varepsilon^\top C\bar\varepsilon)(\bar\varepsilon^\top D\bar\varepsilon)\big],
% \]
% \[
% \overline U := \overline M_{se}\,\overline M_{se}^{\!\top},\qquad
% \overline V := \overline M_{se}\,\overline\Sigma\,\overline M_{se}^{\!\top},\qquad
% \overline W := \overline M_{se}\,\overline\Sigma^{2}\,\overline M_{se}^{\!\top}.
% \]
% \end{proof}



% \section{Proof of Theorem \ref{thm:lqr-lifted-state-maps} }
% \label{sec:app-proof-lqr-lifted-state-maps}

% \begin{proof}
    
% From \eqref{eq:lqr-lift-state-maps}, we have the following summation formulas:
% \begin{equation}\label{eq:FS-gramians}
% \mathcal F_S^\top\mathcal F_S=\sum_{t=0}^{T-1}(F^\top)^t F^t,
% \qquad
% \mathcal F_S^{+\top}\mathcal F_S^+=\sum_{t=1}^{T}(F^\top)^t F^t.
% \end{equation}
% Consequently, for the spectral norm $\|\cdot\|_2$,
% \begin{align}
% \max_{0\le t\le T-1}\|F^t\|_2^2
% \;\le\;
% \|\mathcal F_S\|_2^2
% \;\le\;
% \sum_{t=0}^{T-1}\|F^t\|_2^2, \label{eq:FS-sandwich}\\[4pt]
% \max_{1\le t\le T}\|F^t\|_2^2
% \;\le\;
% \|\mathcal F_S^+\|_2^2
% \;\le\;
% \sum_{t=1}^{T}\|F^t\|_2^2. \label{eq:FSp-sandwich}
% \end{align}

% Each summand $(F^t)^\top F^t\succeq 0$, so by Loewner order
% \(
% (F^{t_0})^\top F^{t_0}\preceq \sum_{t}(F^t)^\top F^t
% \)
% for any $t_0$, which gives the lower bounds in
% \eqref{eq:FS-sandwich}–\eqref{eq:FSp-sandwich} after taking spectral norms.
% For the upper bounds, use the triangle inequality and submultiplicativity:
% \[
% \Big\|\sum_t (F^t)^\top F^t\Big\|_2
% \le \sum_t \|(F^t)^\top F^t\|_2
% \le \sum_t \|F^t\|_2^2.
% \]
% Thus, $\|\mathcal F_S\|_2^2$ is of order $O(T)$ if $\rho(F)=1$; of order $O(1)$ if $\rho(F)<1$; and of order $O(\rho(F)^{2T})$ if $\rho(F)>1$.

% Similar analysis applies to $\mathcal F_E$ and $\mathcal F_E^+$: Note that $ \|\mathcal F_E\|_2 \leq  \|\mathcal F_E^+\|_2$ because $\mathcal F_E$ is just a "shift-down" of $\mathcal F_E^+$. $||\mathcal F_E v|| = ||\mathcal F_E^+ \mathcal S v||$, here $\mathcal S$ is a "shift-down" operator, thus by definition of operator norm, we have $ \|\mathcal F_E\|_2 \leq  \|\mathcal F_E^+\|_2$.

% For $k\ge 1$ define the partial observability Gramian
% \[
% W_k := \sum_{s=0}^{k-1}(F^\T)^s F^s \in \mathbb R^{n\times n}.
% \]
% By definition,
% \[
% \big[(\calF_E^+)^\T \calF_E^+\big]_{i,j}
% = \sum_{t=1}^{T} \big([\calF_E^+]_{t,i}\big)^\T \,[\calF_E^+]_{t,j}
% = \sum_{t=\max\{i,j\}}^{T} (F^{\,t-i}B)^\T (F^{\,t-j}B).
% \]
% For $i\le j$,
% \[
% \sum_{t=j}^{T} B^\T (F^\T)^{\,t-i} F^{\,t-j} B
% = B^\T (F^\T)^{\,j-i} \Big(\sum_{s=0}^{T-j} (F^\T)^{\,s} F^{\,s}\Big) B
% = B^\T (F^\T)^{\,j-i} W_{T-j+1} B,
% \]
% which is \eqref{eq:block-formula}. The $j<i$ case follows by symmetry.
% Then for $1\le i\le j\le T$,
% \begin{equation}\label{eq:block-formula}
% \big[(\calF_E^+)^\T \calF_E^+\big]_{i,j} \;=\;
% B^\T (F^\T)^{\,j-i}\, W_{T-j+1}\, B,
% \qquad
% \big[(\calF_E^+)^\T \calF_E^+\big]_{j,i} \;=\; \big(\big[(\calF_E^+)^\T \calF_E^+\big]_{i,j}\big)^\T .
% \end{equation}

% \end{proof}